%======================================================================
%	Vorlage
%======================================================================
%	$Id$
%	Matthias Kupfer
%======================================================================
%	Documentclass
%======================================================================
\documentclass[
%	draft,			% Entwurfsmodus: Bilder als Rahmen,
				% berlngen werden deutlich markiert
%	10pt,			%
	11pt,			% KOMA default
%	12pt,			%
	a4paper,		% DIN A4
	twoside,		% Zweiseitig
	german,			%
%----------------------------------------------------------------------
% Die folgenden Befehle stammen aus dem KOMA Paket
	headsepline,		% Linie unter der Kopfzeile
%	headnosepline,		%
%	foodsepline,		% Linie ber Fussnote
	footnosepline,		%
	automark,		% Kolumnentitel lebendig
%	bigheadings,		% default: berschriften gross setzen
	normalheadings,		% berschriften normal setzen
%	smallheadings,		% berschriften eher klein setzen
%	pointlessnumbers,	% Keinen Punkt hinter die letzte Zahl
				% eines Kapitels (auch bei Anhang)
%	chapterprefix,		% Kapiteluberschriften "Kapitel"
	appendixprefix,		% Anhang
%	openright,		% Kapitel auf der rechten (ungeraden) Seite
				% anfangen lassen
	openany,		%
%	cleardoublestandard,	%
	cleardoubleplain,	%
%	cleardoubleempty,	%
	abstracton,		%
	idxtotoc,		% Index soll im Inhaltsverzeichnis auftauchen
	liststotoc,		%
	bibtotoc,		%
%	parskip, 		% parskip-, parskip*, parskip+
% 	halfparskip, 		% halfparskip-, halfparskip* und halfparskip+
%	DIVclassic,
 	BCOR8mm,		%%?? Bindungskorrektur:
				% BCOR<Breite des Bindungsverlustes>
]{scrreprt}

%======================================================================
%	Bearbeitung sowohl mit LaTeX als auch mit pdfLaTeX ermoeglichen
%======================================================================
%\newif\ifpdf
	%\ifx\pdfoutput\undefined
	%\pdffalse 			% we are not running pdflatex
%\else
	%\pdfoutput=1 			% we are running pdflatex
	%\pdftrue
%\fi


%======================================================================
%	Verwendete Pakete
%======================================================================

%\usepackage{babel}		% Sprachen

%% Latex mit deutschen Umlauten:
%% http://www.cs.albany.edu/~herrmann/latex_umlaute/
\usepackage[utf8]{inputenc}   % Eingabe von ,,, erlaubt

\usepackage[T1]{fontenc}	% EC-Schriften verwenden (vs. DC) da 8-Bit
				% EC-Schriften als T1-kodierten CM-Schriften
				% European/Ext.-Computer-Modern-(EC)-Schriften
				% Umlaute, Anfhrungszeichen ...
				% => Umlauten koennen richtig getrennt werden
				% FAQ 5.3.2

\usepackage{ae,aecompl}		% virtuelle-CM-Fonts
				% da EC nicht als PostScript-(Type-1) verfuegbar
				% => keine echten Umlaute im PDF-Dokumen
				%(Problem bei Suche)
				% By loading the ae package (\usepackage{ae}),
				% you loose some characters as mentioned in
				% README.
				% The package aecompl by Denis Roegel restores
				% these characters which are taken from the ec
				% fonts. If you use pdftex, you will get these
				% characters as bitmaps, but this might be
				% better than not having them at all.

%\usepackage{times, mathptm}	% TimesNewRoman Schrift (Acrobat Reader Fonts),
				% dazu braucht man auch den entsprechenden
				% Zeichensatz fr den Math-Mode
%\usepackage{pslatex}		% ? mathematische Formeln mit Standard
				% Postscript Fonts gesetzt
			% Paket     Roman         Serifenlos  Typewriter
%\usepackage{times}	% -----------------------------------------------
			% times     Times         Helvetica   Courier
			% palatino  Palatino      Helvetica   Courier
			% newcent   NewCenturySch AvantGarde  Courier
			% bookman   Bookman       AvantGarde  Courier
			% Diese Schriften sind die Standard-PostScript-Schriften
			% und in jedem Drucker verfgbar

\usepackage{
%	german,			% Deutsche Trennungen (ALTE Rechtschreibung),
				% Anfhrungsstriche und mehr
	ngerman,		% Deutsche Trennungen (NEUE Rechtschreibung),
				% Anfhrungsstriche und mehr
				%
%	acronym,		% Verwaltung von Abkuerzungen
	bibgerm,		% Deutsche Bibliographie
	calc,			% Erweiterung der arithmetischen Funktionen in
				% LaTeX
				% wird verwendet um Titelseite zu zentrieren
	color,			% im Laufenden Text einfach mit \color{Farbe) zwischen den
				% Farben umschalten, wobei Farbe einfach
				% durch z.B. red, blue, black etc. ersetzt wird
				% \textcolor{farbe){Text)
%	epigraph,		% Zitat am Kapitelanfang
%	fancyhdr,		% Kopf- und Fuzeilen von Dokumenten frei
				% gestalten
	fancybox,		% shadowbox, doublebox, ovalbox, Ovalbox
	fancyvrb,		% verbatim Erweiterung:
	float,			% Positionierung von Gleitobjekten genau an der Stelle, wo man
				% 'figure'- oder 'table'-Umgebung die
				% Positionierung [H] gesetzt werden
%	glosstex,		% Glossar und Abkrzungsverzeichnis
	mdwlist,		% compact list: itemize* ..
	scrdate,		% \todaysname
	scrtime,		% \thistime
	scrpage2,		% Kopf- und Fuzeilen flexibel gestalten
				%
%	moreverb,		% verbatim-hnlich: boxedverbatim, listing
%	verbatim,		% Darstellung von "Text, wie er eingegeben wird"
				%
%	lscape,			% Erstellt eine um 90% gedrehte *neue* Seite
%	textcomp,		% Sonderzeichen
%	booktabs,		% Tabellenlinien
%	longtable,		% Tabellen > 1 Seite
%	supertabular,		% Tabellen > 1 Seite
	tabularx,		% Blocksatzspalten
%	ltxtable,		% tabularx + longtable
%	multicol,		% mehrspaltige Zeilen
%	varioref,		% einheitliche Verweise
%	endnotes,		% Fussnoten -> Endnoten
%	rotating,		% sidewaystable und sidewaysfigure
%	natbib,			% Bibliographie ohne Klammer etc.
%	marvosym,		% Euro etc.
}

%\usepackage{pstricks}
%\usepackage{listings}
%\usepackage{wasysym}
%\usepackage{caption}
%\usepackage{subcaption}
\usepackage{subfig}
%\usepackage[german, first, bottomafter]{draftcopy}

%\usepackage{setspace}	% Durchschu, Zeilenabstand
%\doublespace		% doppelzeilig oder
%\onehalfspacing	% anderthalbzeilig

% Fr schne Darstellung von Algorithmen
%\usepackage[german, algoruled, algochapter]{algorithm2e}
%\usepackage{algorithmic}
%\usepackage[chapter]{algorithm}
%\floatname{algorithm}{Algorithmus}

%======================================================================
%	Bilder, Links
%======================================================================

\usepackage{graphicx}
%\graphicspath{{../img/}}	% Angabe der Pfade, wo die Grafiken liegen;
				% mehrere Pfade sind mglich

\usepackage{thumbpdf}
\usepackage{hyperref}


%======================================================================
%	user.sty
%======================================================================

\usepackage{user}	% Makros in user.sty
			% \epsinc{bild.eps}{scale=1}{Bildunterschrift}
			% \missing{Beweis fehlt noch}
			% \comment{Ein Kommentar}

%======================================================================
%	Einstellungen
%======================================================================

%\typearea{11}		% Satzspiegel neu konstruieren (KOMA)
			%      10pt 11pt 12pt
			% DIV : 8   10   12

\pagestyle{scrheadings}	% Standart  Kopf- und Fuzeile
\setkomafont{pagehead}{\small\scshape}

% ---------------------------------------------------------------------
%\setcounter{secnumdepth}{2}
%\setcounter{chapter}{-1}
\setcounter{tocdepth}{3}

% ---------------------------------------------------------------------
%\sloppy		% weniger Worttrennungen, grere Wortabstnde
\fussy			% viele Worttrennungen, "schnere" Wortabstnde

% ---------------------------------------------------------------------
%\flushbottom			% Ausrichtung der Seitenenden jeweils auf
				% gleicher Hhe

% ---------------------------------------------------------------------
%\sloppypar		% Das hier relaxt die Einstellungen zum Wortabstand
			% extrem. Damit ragen keine Worte ber den rechten
			% Zeilenabstand hinaus. Dafr mu strker auf
			% Wortabstand geachtet werden, der kann dann ziemlich
			% gro werden. Man erhlt aber keine Meldung mehr
			% ber underfull boxes.

%Hiermit kann man das gleiche mit weniger Holzhammer erreichen:
%\setlength{\tolerance}{2000}           % Strafpunkt fr Zeilenumbruch
%\setlength{\emergencystretch}{3pt}     % Soweit drfen einzelne Worte mehr
					% auseinandergezogen werden
%\setlength{\hfuzz}{1pt}                % Macht den rechten Rand um bis zu 1pt
					% flatterig.

% ---------------------------------------------------------------------
%Vermeiden einzelner Zeilen am Ende einer Seite oder oben auf einer neuen Seite
%\clubpenalty10000
%\widowpenalty10000

% ----------------------------------------------------------------------
% neue Umgebungen fr verwendete Stze und Beispiele
\newtheorem{bsp}{Beispiel}[chapter]
\newtheorem{satz}{Satz}[chapter]

%======================================================================
%	includeonly
%======================================================================

\includeonly{		% Gibt an, welche Dateien der include-Befehl
			% tatschlich einfuegen darf.
  metadaten		% Variablen setzen
  ,titel		% Titelseite, Zusammenfassung und Inhaltsverzeichnis
% Alle per include einzulesenden Dateien mssen hier angegeben sein!
  %,anleitung		% Anleitung zur Nutzung der Vorlage
	,main
}

% Aufteilung der Arbeit in folgende Bestandteile in angeg. Reihenfolge
%
% Titelseite
% Bibliographische Beschreibung (Rckseite der Titelseite)
% (Aufgabenstellung)
% Danksagung
% Abstract
% Inhaltsverzeichnis
% Tabellenverzeichnis
% Abbildungsverzeichnis
% Einleitung
% Inhalt
% Zusammenfassung/Ausblick
% (Verzeichnis verwendeter Terme)
% Glossar (Abkrzungsverzeichnis)
% (Index)
% Literaturverzeichnis
% (Thesen)
% (Selbststndigkeitserklrung)
% Anhnge


%======================================================================
% * *   T E X T    S T A R T S    H E R E   * * * * * * * * * * * * * *
%======================================================================

\begin{document}
%======================================================================
%	Metadaten
%======================================================================
%	$Id$
%	Matthias Kupfer
%======================================================================

\newcommand{\dcsubject}{Praktikumsbericht}
% z.B. (Diplom/Studien/Haus)arbeit, Praktikumsbericht, Studie, Beleg,
% (Pro/Haupt/Ober)seminar, Seminar usw.
\newcommand{\dctitle}{MobileMixedReality Report}
\newcommand{\dcsubtitle}{~} % Untertitel, falls erforderlich

\newcommand{\dcauthorlastname}{David Eilers, Oskar Hippe}
\newcommand{\dcauthorfirstname}{}
\newcommand{\dcauthoremail}{name@subdomain.domain}
\newcommand{\dcdate}{\today}

\newcommand{\dcplace}{Chemnitz} % Ort, kann an der TU meist so bleiben
\newcommand{\dcuni}{Technische Universitt \dcplace}
\newcommand{\dcdepart}{Fakultät Informatik} % Fakulttsangabe
\newcommand{\dcprof}{Professur Graphische Datenverarbeitung und Visualisierung} % Angabe der Professur

\newcommand{\dcpruefer}{Prof. Dr. Guido Brunett}% Prfer der Arbeit
\newcommand{\dcadvisor}{Tom Uhlmann M.Sc.}% Betreuer der Arbeit

\newcommand{\dckeywords}{Liste,von,Stichworten,die,als,Schlagworte,geeignet,
sind}

%%======================================================================
% Einstellungen des Hyperref-Paketes
\hypersetup{%
	pdftitle	= {\dctitle}, %
	pdfsubject	= {\dcsubject, \dcdate}, %
	pdfauthor	= {\dcauthorfirstname~\dcauthorlastname, \dcauthoremail}, %
	pdfkeywords	= {\dckeywords}, %
	pdfcreator	= {pdfTeX with Hyperref and Thumbpdf}, %
	pdfproducer	= {LaTeX, hyperref, thumbpdf}, %
	% weitere PDF-Einstellungen in hyperref.cfg
}
%%======================================================================
 % Variablen, \hypersetup, etc.

%======================================================================
%	Titelseite (Zusammenfassung und Inhaltsverzeichnis)
%======================================================================

%======================================================================
%	Titelseite
%======================================================================
%	$Date:$
%	$Revision:$
%	Matthias Kupfer
%======================================================================

%%======================================================================
%% Schmutztitel
%%======================================================================
%\extratitle{
%	\usekomafont{sectioning}\mdseries
%	\begin{center}
%		\Huge \dcsubject\\[1.5ex]
%		\hrule
%		\vspace*{\fill}
%		\includegraphics{TUC_deutsch_einzeile_CMYK}
%	\end{center}
%}

%%======================================================================
%% Titelkopf
%%======================================================================
\titlehead{
	\vspace*{-1.5cm}
	% Schriftfamilie wie alle berschriften, aber nicht fett
	\usekomafont{disposition}\mdseries
	\begin{center}
		\raisebox{-1ex}{\includegraphics[scale=1.4]{TUC_deutsch_einzeile_CMYK}}\\
		\hrulefill \\[1em]
		{\Large\dcdepart}\\[0.5em]
		\dcprof
	\end{center}
	\vspace*{1.5cm}
}

%%======================================================================
%% Subjekt
%%======================================================================
\subject{\bf\Huge\dcsubject}


%%======================================================================
%% Titel
%%======================================================================
\title{\sf\Large
	\dctitle
	\\
	\dcsubtitle
}

%%======================================================================
%% Autor des Dokumentes
%%======================================================================
\author{\dcauthorfirstname~\dcauthorlastname}

%%======================================================================
%% Ort, Datum
%%======================================================================
\date{\dcplace, den \dcdate
}

%%======================================================================
%% Publishers
%%======================================================================
\publishers{
	{\parbox{\textwidth-8em}{
		\begin{tabbing}
			{\bf Betreuer:}\quad\=\kill
			{\bf Prüfer:}	\>\dcpruefer\\
			{\bf Betreuer:}	\>\dcadvisor
		\end{tabbing}
	}}
}

%%======================================================================
%% bibliografische Angaben
%%======================================================================
\lowertitleback{
\textbf{\dcauthorlastname, \dcauthorfirstname}\\
\dctitle\\
\dcsubject,~\dcdepart\\
\dcuni,~\ifcase\month\or
  Januar\or Februar\or März\or April\or Mai\or Juni\or
    Juli\or August\or September\or Oktober\or November\or Dezember\fi
    ~\number\year
}

%%======================================================================
%% maketitle
%%======================================================================

\maketitle

%%======================================================================
%% Danksagung
%%======================================================================
% \thispagestyle{empty}
% \null\vfil
% \begin{center}
% \usekomafont{disposition}\textbf{Danksagung}
% \vspace{-.5em}\vspace{\parsep}
% %%
% %% Hier steht der Text fr die Danksagung
% %%
% \end{center}
% \par\vfil\null
% \cleardoubleemptypage

%%======================================================================
%%      Kurzfassung / Abstract
%%======================================================================
%%\def\abstractname{Abstract} 	% Wenn der Text "Zusammenfassung" erscheinen
				% soll, dann mu dies auskommentiert werden

%%\begin{abstract}
%%In diesem Praktikum war es unsere Aufgabe eine augmented reality App zu entwickeln,
in der zwei Spieler auf einem Gerät oder auf zwei Geräten im lokalen Netzwerk gegeneinander
Vier gewinnt spielen können. Dabei sind wir bei der Wahl der AR-Umgebung auf mehrere Sackgassen
gestoßen, haben aber mit AR-Core eine ausreichende Lösung gefunden. Die App wurde hauptsächlich in C++ geschrieben.
Um dies auf Android zu ermöglichen nutzten wir das Android NDK (native development kit). Damit ist es möglich
den für Android üblichen Jave-Code über das Java Native Interface (JNI) mit dem C++ Code kommunizieren
zu lassen. Nicht nur der grafische Teil der App sondern auch die Netzwerkfunktionen wurden mit C++ umgesetzt.
Lediglich die Activities und Menüelemente wurden in Java implementiert.
Beim Testen der App stießen wir auf mehrere Probleme. Zum einen sind Smartphones noch nicht flächendeckend
AR-Core-fähig. Ältere Geräte sind dadurch für Tests ausgeschlossen. Zum anderen ist der Umgang mit dem Android-Emulator
noch schwierig und teilweise garnicht möglich, wenn der PC keine Virtualisierung unterstützt.
Zum Testen der Netzwerkfunktionen kam deshalb ein Python-Skript zum Einsatz.
Die Entwicklung der Netzwerkschnittstelle war insgesamt gesehen aber unkompliziert im Vergleich zum grafischen Teil.
Die App war am Ende des Praktikums Funktionsfähig und kann auf AR-Core-fähigen Smartphones ausgeführt werden.
Für die Zukunft wäre es jedoch Vorteilhafter die Anwendung entweder vollständig in Java oder vollständig in C++
zu entwickeln. Das JNI ist kompliziert und langsam und dadurch hätte man auch auf eine native avtiviy setzen können.
Auf diese Weise wäre der eigentliche Code der App vollstädig in C++ gewesen. Die App ist aber auch
problemlos in Java umsetzbar und im Context von Android ist Java als Programmiersprache die unkompliziertere Wahl.
%%\end{abstract}

%%======================================================================
%%      Inhaltsverzeichnis
%%======================================================================
%\cleardoubleemptypage
\pagenumbering{roman}
\pdfbookmark{Inhaltsverzeichnis}{Inhaltsverzeichnis}
\tableofcontents

%%======================================================================
%%      Abbildungsverzeichnis
%%======================================================================
%\cleardoublepage
%\markboth{Abbildungsverzeichnis}{Abbildungsverzeichnis}
%\listoffigures

%%======================================================================
%%      Tabellenverzeichnis
%%======================================================================
%\cleardoublepage
%\markboth{Tabellenverzeichnis}{Tabellenverzeichnis}
%\listoftables

%%======================================================================
%%      Algorithmenverzeichnis
%%======================================================================
%\renewcommand{\listalgorithmname}{Algorithmenverzeichnis}
%\cleardoublepage
%\addcontentsline{toc}{chapter}{Algorithmenverzeichnis}
%\listofalgorithms

%%======================================================================
%%      Abkuerzungsverzeichnis
%%======================================================================
%\cleardoublepage
%\addcontentsline{toc}{chapter}{Abkrzungsverzeichnis}
%\markboth{Abk"urzungsverzeichnis}{Abk"urzungsverzeichnis}
%\def\listacronymname{Abk"urzungsverzeichnis}
%\printglosstex(acr)

%%======================================================================
%%      Ende
%%======================================================================
%%\cleardoublepage
\pagenumbering{arabic}
\setcounter{page}{1}
 % Titelseite, Zusammenfassung
		% Inhaltsverzeichnis
		% Abkuerzungsverzeichnis
		% -> danach Seitennummerierung bei 1

%======================================================================
%	Inhalt per include
%======================================================================

%

\section{Einleitung}
Diese Datei enthlt die Anleitung zur Nutzung der Vorlage fr verschiedene
Typen von Arbeiten. Sie ist vorrangig fr Studenten (und auch wissenschaftliche
Mitarbeiter) gedacht, welche ihre Arbeiten bzw. Publikationen mit \LaTeX{}
erstellen wollen. Dabei wurden auch die Richtlinien des Corporate Design
der Technischen Universitt Chemnitz bercksichtigt, soweit sie sich
ohne grere Probleme in \LaTeX{} realisieren lassen.

Diese Vorlage ist fr folgende Dokumente konzipiert, kann aber bei geringen
Modifikationen auch darber hinaus eingesetzt werden:
\begin{itemize}
\item Hausarbeiten
\item Studienarbeiten
\item Diplomarbeiten
\item Praktikumsberichte
\item Proseminare
\item Oberseminare
\item Hauptseminare
\item sonstige Seminare
\item Belege
\item Studien
\end{itemize}

In den folgenden Kapiteln dieser Anleitung wird ein berblick ber die
Verwendung der Vorlage, Zweck der Dateien und typische Anwendungsflle
gegeben.

\textbf{Hinweis:} Diese Anleitung ist \textbf{keine} Einfhrung
in \LaTeX. Dazu sei auf das Kursangebot des URZ bzw. auf weiterfhrende
Literatur verwiesen. Auch erhebt diese Vorlage \textbf{nicht} den Anspruch,
da jedes damit erstellte Dokument innerhalb der TU-Chemnitz grundstzlich
in Form, Umfang und Aufbau anerkannt wird. Studenten sollten dies
grundstzlich vor der Verwendung anhand der fr sie gltigen Studien-
und Prfungsordnung prfen und darber hinaus mit dem fr sie
zustndigen Professor bzw. Betreuer klren.
\section{Grundkonzept}
Alle auf Basis dieser Vorlage erstellten Dokumente verwenden das sogenannte
{\scshape Koma-Script}-Paket. Dieses Paket wurde entwickelt um \LaTeX{} den
europischen Anforderungen (insbesondere Deutschland) anzupassen.
Diese Anforderungen umfassen u.a.
\begin{itemize}
\item Papierformate
\item verschiedene Sprachen
\item verschiedene Datumsformate
\end{itemize}
Fr weitere Details sei auf die Dokumentation zu {\scshape Koma-Script}
verwiesen.
 % Datei, welche die Anleitung zur Nutzung der
		    % Vorlage enthlt
\chapter{Introduction}
\section{Zusammenfassung}
Dies ist der Praktikumsberichte zu unserem Teampraktikum. In diesem haben wir
eine AR App entwickelt, in der man das Spiel ´´Vier gewinnt´´ spielen kann. 




%======================================================================
%	Glossar
%======================================================================

%\manualmark
%\addcontentsline{toc}{chapter}{Glossar}
%\markboth{Glossar}{Glossar}
%\def\glossaryname{Glossar}
%\printglosstex(glo)
%\cleardoublepage

%======================================================================
%	Literaturverzeichnis
%======================================================================
\manualmark
\markboth{Literaturverzeichnis}{Literaturverzeichnis}
\bibliographystyle{user}
\bibliography{literatur}
%\cleardoublepage

%======================================================================
%	Thesen
%======================================================================
%\chapter{Thesen}
%\include{thesen}

%======================================================================
%	Selbststndigkeitserklrung
%======================================================================
\chapter*{Selbstständigkeitserklärung}

Hiermit erkläre ich, da ich die vorliegende Arbeit
selbstständig angefertigt, nicht anderweitig zu Prüfungszwecken vorgelegt und
keine anderen als die angegebenen Hilfsmittel verwendet habe. Sämtliche
wissentlich verwendete Textausschnitte, Zitate oder Inhalte anderer Verfasser
wurden ausdrcklich als solche gekennzeichnet.\\[2ex]
\dcplace, den \dcdate\\[6ex]
\flushleft
\newlength\us
\settowidth{\us}{-\dcauthorfirstname~\dcauthorlastname-}
\begin{tabular}{p{\us}}\hline
\centering\footnotesize \dcauthorfirstname~\dcauthorlastname
\end{tabular}

%======================================================================
%	Anhang
%======================================================================


%\part*{Anhang}
%\cleardoublepage
\appendix
%\include{anhanga}

\end{document}
